
\[
\begin{tikzcd}[column sep=large,row sep=large]
B\otimes_A B \arrow[r,"p\otimes p"] \arrow[d,"\Delta"'] \arrow[dr,phantom,"(1)"] & B'\otimes_A B' \arrow[d,"\Delta"]\\
B \arrow[r,"p"] \arrow[dr,"i"'] & B' \arrow[d,"j"'] \arrow[dr,"i'","(2)"' near start] & B'\otimes_B B' \arrow[l,"\Delta"'] \arrow[d]\\
{} & C \arrow[r,shift right=.7ex,"p'"'] & B'\otimes_B C = B'/IB'. \arrow[l,shift right=.7ex,"j'"']
\end{tikzcd}
\]

L'homomorphisme $\underline{I}/\underline{I}^2\to\underline{J}/\underline{J}^2$ se déduit de l'homomorphisme : $IB'\to J$ et coïncide avec celui qui se déduit de $p:I\to J$; l'homomorphisme : $i^*(\underline{\Omega}^1_{V/Y})\longrightarrow j^*(\underline{\Omega}^1_{V'/Y})$ se déduit du carré commutatif (1). La propriété a) est donc évidente.

L'homomorphisme : $\underline{K}/\underline{K}^2\to j^*(\underline{\Omega}^1_{V'/V})$ considéré dans b) peut être défini comme suit: soit $\overline x\in K/K^2$, soit $x$ un élément de $J$ remontant $\overline x$ (il en existe car $K=J/IB'$). On associe à $\overline x$ l'image de $x$ par l'homomorphisme composé : $J\xrightarrow{d}\Omega_{B'/A}\otimes_{B'}C\longrightarrow\Omega^1_{B'/B}\otimes_{B'}C$. Par ailleurs, soit $H$ l'idéal d'augmentation de : $B'\otimes_B B'\longrightarrow B'$ ; le carré commutatif (2) donne un homomorphisme : $H\to K$, d'où un homomorphisme : $H/H^2\otimes_{B'}C\longrightarrow K/K^2$. Or $H/H^2=\Omega^1_{B'/B}$ ; par suite, il suffit de vérifier que les homomorphismes construits sont inverses l'un de l'autre, ce qui se fait sans difficulté.

\textbf{Définition 2.3.} \uline{Soit} $f:X\to Y$ \uline{un morphisme lissifiable}. \uline{Le complexe} $\underline{L}^{X/Y}_\bullet$, \uline{défini à isomorphisme canonique près dans la catégorie dérivée} $D(X)$, \uline{est appelé} complexe cotangent relatif \uline{de} $X$ \uline{sur} $Y$ (\uline{ou de} $f$).

\textbf{Proposition 2.4.} \uline{Soit} $f:X\to Y$ \uline{un morphisme d'intersection complète} (1.1) \uline{lissifiable}. \uline{Alors} $\underline{L}^{X/Y}_\bullet$ \uline{est parfait}. \uline{Plus précisément}, \uline{pour toute}
